
Per-benchmark threshold calibration was unexpected---initial planning specified uniform 0.5\% based on ImageNet historical patterns. Synthetic data artifact is the primary explanation, with potential domain factor (vision vs. NLP variance differences) as secondary contributor. Real PWC data will determine whether calibration is genuinely required or a synthetic-only issue.

Citation correlation precision failure suggests the 6-month correlation window may be too narrow for real-world paradigm shift adoption patterns. GLUE false positive (7-month shift, only 1 month outside threshold) indicates edge case sensitivity. Expanding to n = 15+ benchmark-shift pairs and testing window variants (9-month alignment) are necessary next steps.


Machine learning benchmarks drive substantial investment without systematic saturation detection mechanisms. Temporal precedence analysis reveals that saturation signals appear an average of 4 years before community migration---far longer than previously assumed, providing an actionable early warning window for proactive benchmark rotation.

Through proof-of-concept experiments on major benchmarks (ImageNet, GLUE, SQuAD), three key contributions emerged. First, expert consensus on saturation timing exists with 76--93\% agreement within $\pm 1$ year of modal dates, validating saturation as a community-observable phenomenon. Second, score convergence detection via rolling window statistics achieves 100\% detection rate with statistical significance (Levene's p < 0.05), demonstrating algorithmic detectability without complex models. Third, all tested saturations preceded paradigm shift adoption by 32--78 months (mean 48 months), distinguishing internal benchmark exhaustion from external disruption.

These findings provide validated mechanisms that, pending real-world deployment, could enable a shift from reactive organic migration to systematic benchmark lifecycle management. Treating benchmarks as consumables with time-boxed validity periods provides the foundation for saturation detection infrastructure. Conference organizers can integrate saturation dashboards into submission systems, enabling benchmark selection rationale disclosure and voluntary adoption.
